\renewcommand{\Anno}{2012}

\renewcommand{\NumeroProblema}{3}

\renewcommand{\Testo}{ In una strada rettilinea ci sono $n$ lampioni.
Uno è posto all'inizio, un altro alla fine, e la distanza tra due lampioni consecutivi è sempre di 10 metri.  Ogni lampione illumina un cerchio con raggio di 12 metri centrato alla sua base.

In una notte, ogni lampione ha una probabilità $p$ di spegnersi. Indichiamo con $P_{n}(p)$ la probabilità che al termine della notte esista almeno un punto della strada non illuminato da lampioni (si trascuri l'ampiezza della carreggiata e la distanza dei lampioni dalla stessa).
% {\cyan [Ogni altra soluzione mi sembra peggiore del
% male: non e' di certo credibile che la zona illuminata sia un
% rettangolo!  e poi c'e' il problema del lampioni ai lati che non 
% illuminano in entrambe le direzioni \ldots non vorrei far venire dei 
% dubbi a chi non li aveva \ldots]}

\begin{enumerate}
	\renewcommand{\labelenumi}{(\alph{enumi})}
	\item Determinare $P_{4}(p)$.

	\item  Determinare una formula che permetta di calcolare $P_{n+1}(p)$ in funzione di $P_{n}(p)$ e $P_{n-1}(p)$.
\end{enumerate}

% {\cyan [Ho scelto la strada con il massimo del suggerimento. Diversamente mi 
% sembra un esercizio che rischia di andare deserto o riservato ai 
% ``professionisti'' (che poi non prendono il posto).]}
}

\renewcommand{\Risposta}{
\begin{enumerate}
	\renewcommand{\labelenumi}{(\alph{enumi})}
	\item $P_{4}(p)=3p^{2}-2p^{3}$.

	\item  $P_{n+1}(p)=p^{2}+(1-p)P_{n}(p)+p(1-p)P_{n-1}(p)$.
\end{enumerate}
}

\renewcommand{\Aiutino}{
\begin{enumerate}
	\renewcommand{\labelenumi}{(\alph{enumi})}
	\item Quali sono le possibili configurazioni? Qual è la probabilità di ognuna di esse?

	\item Consideriamo una configurazione con $(n+1)$ lampioni.  Se l'ultimo è acceso, decidono i precedenti, quindi \ldots\ Se invece l'ultimo è spento, si guarda il penultimo: se è spento, allora \ldots; se è acceso, allora \ldots
\end{enumerate}
}

\renewcommand{\SoluzioneUfficiale}{\TitoloSoluzione
Osserviamo preliminarmente che al termine della notte ci sarà un punto della strada al buio se e solo se si saranno spenti almeno 2 lampioni consecutivi. Possiamo dunque pensare $P_{n}(p)$ come la probabilità che almeno 2 lampioni consecutivi siano spenti.

Mostriamo direttamente che
\begin{equation}
	P_{n+1}(p)=p^{2}+(1-p)P_{n}(p)+p(1-p)P_{n-1}(p),
	\label{ricorrenza-Pn(p)}
\end{equation}
il che risponde al punto~(b). 

Poiché ovviamente $P_{1}(p)=0$ (se c'è un solo lampione, non ce ne possono essere 2 consecutivi spenti) e $P_{2}(p)=p^{2}$ (se ci sono 2 lampioni, questi devono essere entrambi spenti), utilizzando
la (\ref{ricorrenza-Pn(p)}) ricaviamo facilmente che
$$P_{3}(p)=p^{2}+(1-p)P_{2}(p)+p(1-p)P_{1}(p)=p^{2}+(1-p)p^{2}=2p^{2}-p^{3}$$
e
\begin{eqnarray*}
	P_{4}(p) & = & p^{2}+(1-p)P_{3}(p)+p(1-p)P_{2}(p)  \\
	 & = & p^{2}+(1-p)(2p^{2}-p^{3}) +p(1-p)p^{2} \\
	 & = & 3p^{2}-2p^{3},
\end{eqnarray*}
il che risponde al punto~(a).

Per dimostrare la (\ref{ricorrenza-Pn(p)}), prendiamo una 
configurazione con $(n+1)$ lampioni, e fissiamo l'attenzione sul lampione $(n+1)$-esimo.
\begin{itemize}
	\item Se il lampione $(n+1)$-esimo è acceso, cosa che accade con probabilità $(1-p)$, avremo che i 2 lampioni consecutivi spenti dovranno eventualmente stare tra i primi $n$, cosa che accade con probabilità $P_{n}(p)$.  Mettendo insieme le 2 cose, questo caso contribuisce con $(1-p)P_{n}(p)$ al calcolo di $P_{n+1}(p)$.

	\item  Se il lampione $(n+1)$-esimo è spento, cosa che accade con probabilità $p$, spostiamo l'attenzione sull'$n$-esimo, il che ci porta a distinguere 2 sottocasi.
	\begin{itemize}
		\item Se anche il lampione $n$-esimo è spento, allora abbiamo già trovato 2 lampioni consecutivi spenti, e non dobbiamo cercare oltre.  Questo sottocaso contribuisce con $p^{2}$ (cioè la probabilità che siano spenti sia l'$n$-esimo, sia l'$(n+1)$-esimo) al calcolo di $P_{n+1}(p)$.
	
		\item Se il lampione $n$-esimo è acceso, allora i 2 lampioni consecutivi spenti dovranno eventualmente stare tra i primi $(n-1)$, cosa che accade con probabilità $P_{n-1}(p)$.  Pertanto questo ultimo sottocaso contribuisce con $p(1-p)P_{n-1}(p)$ al calcolo di $P_{n+1}(p)$.  Il fattore $p$ rappresenta la probabilità che l'$(n+1)$-esimo lampione sia spento, il fattore $(1-p)$ la probabilità che l'$n$-esimo lampione sia acceso, il fattore $P_{n-1}(p)$ la probabilità che almeno 2 lampioni consecutivi tra i primi $(n-1)$ siano spenti.
	\end{itemize}
\end{itemize}

Poiché i 3 casi o sottocasi considerati sono disgiunti ed esauriscono tutte le possibilità, non resta che sommare i 3
contributi per ottenere la (\ref{ricorrenza-Pn(p)}).
}

\renewcommand{\SoluzioniAlternative}{
È possibile rispondere al punto~(a) con un argomento diretto che
non passa per la ricorrenza~(\ref{ricorrenza-Pn(p)}). Essendo $n=4$, 
ci basta infatti considerare separatamente i seguenti pochi casi disgiunti.
\begin{itemize}
	\item  Se tutti e 4 i lampioni sono spenti, allora sicuramente c'è un punto della strada al buio. La probabilità che questo accada è chiaramente $p^{4}$.

	\item  Se esattamente 3 lampioni sono spenti, sicuramente almeno 2 consecutivi sono spenti, quindi c'è comunque un punto della strada al buio. La probabilità di questo caso è $4p^{3}(1-p)$, in cui il fattore $p^{3}$ rappresenta la probabilità che 3 lampioni ben determinati siano spenti, $(1-p)$ la probabilità che il restante lampione sia acceso, 4 il numero di configurazioni con 3 lampioni spenti e uno acceso.

	\item Se esattamente 2 lampioni sono spenti, allora questi devono essere consecutivi.  La probabilità di questo caso è $3p^{2}(1-p)^{2}$, in cui il fattore $p^{2}$ rappresenta la probabilità che 2 lampioni ben determinati siano spenti, $(1-p)^{2}$ la probabilità che i restanti 2 lampioni siano accesi, 3 il numero di configurazioni con 2 lampioni spenti consecutivi su 4.

	\item  Se i lampioni spenti sono 0 o 1, allora non ci sono punti della strada al buio, dunque questo caso non contribuisce al calcolo di $P_{4}(p)$.
\end{itemize}

Mettendo insieme il contributo dei primi 3 casi, e svolgendo i calcoli, otteniamo che
$$P_{4}(p)=p^{4}+4p^{3}(1-p)+3p^{2}(1-p)^{2}=3p^{2}-2p^{3},$$
il che risponde direttamente al punto~(a).
}

\renewcommand{\Osservazioni}{
Grazie alla teoria generale delle successioni per ricorrenza lineari di ordine 2, è possibile passare dalla relazione ricorrente
(\ref{ricorrenza-Pn(p)}) ad una formula esplicita per $P_{n}(p)$ che 
permette di calcolarne il valore per un certo $n$ senza bisogno di calcolare tutti i valori precedenti. La formula che si ottiene è
$$P_{n}(p)=1+\frac{1+p-\Delta}{2\Delta}
\left[\frac{1-p+\Delta}{2}\right]^{n}-\frac{1+p+\Delta}{2\Delta}
\left[\frac{1-p-\Delta}{2}\right]^{n},$$
dove $\Delta=\sqrt{1+2p-3p^{2}}$.

Una volta scritta, la formula può essere verificata elementarmente per induzione.  Per ricavarla, la teoria generale indica che, nel caso
di una ricorrenza del tipo (\ref{ricorrenza-Pn(p)}), la formula per il
termine generico è del tipo
$$P_{n}=aR_{1}^{n}+bR_{2}^{n}+c,$$
dove $R_{1}$ ed $R_{2}$ sono le radici del polinomio
$$x^{2}-(1-p)x-p(1-p),$$
detto polinomio caratteristico della (\ref{ricorrenza-Pn(p)}), ed $a$,
$b$ e $c$ sono opportuni coefficienti (in questo caso sia le 2 radici, sia i 3 coefficienti sono funzioni di $p$).

Vale la pena notare l'analogia formale con la formula per i numeri di Fibonacci, anch'essi definiti mediante una ricorrenza lineare di ordine 2.  Da notare infine che, nonostante l'abbondante presenza di radicali, la formula deve fornire sempre dei polinomi nella variabile $p$, come si può verificare per induzione partendo dalla
solita~(\ref{ricorrenza-Pn(p)}).  }

\renewcommand{\Commenti}{
}
